\section{Consentimiento informado}

Este formato acompaña a los instrumentos de la sección 4 y se presenta a cada participante
antes de iniciar la encuesta o la entrevista. Recoge los elementos que la literatura de
investigación de usuario considera indispensables: explicar qué se va a hacer, qué se pide del
participante, con qué propósito, cómo se usará el material recopilado, con quién se comparte,
durante cuánto tiempo se conserva y el derecho a retirarse en cualquier momento
(de Voil, 2020).

\subsection{Texto del consentimiento}

\begin{uxdestacado}
\textbf{\textcolor{uxnavy}{Propósito.}} Participas en un estudio académico del curso TC4032
Experiencia del usuario y diseño de interfaces, de la Maestría en Inteligencia Artificial
Aplicada del Instituto Tecnológico y de Estudios Superiores de Monterrey. El objetivo es
comprender cómo las personas localizan, validan y utilizan información financiera en su
trabajo, para orientar el diseño de un portal centralizado de datos.

\textbf{\textcolor{uxnavy}{Qué se te pide.}} Responder una encuesta de entre cinco y siete
minutos, o participar en una conversación de entre treinta y cuarenta minutos sobre tu
experiencia de trabajo con datos. No hay respuestas correctas ni incorrectas y no se evalúa tu
desempeño.

\textbf{\textcolor{uxnavy}{Qué no se te pide.}} No solicitamos cifras, nombres de clientes,
credenciales ni información confidencial de tu institución. Si una pregunta te lleva a ese
terreno, omítela.

\textbf{\textcolor{uxnavy}{Uso del material.}} Tus respuestas se emplean únicamente con fines
académicos, dentro de los entregables del curso. Se reportan de forma agregada o anonimizada:
no se publican tu nombre, tu puesto ni el nombre de tu institución. Las citas textuales se
usan sin atribución nominal.

\textbf{\textcolor{uxnavy}{Quién tiene acceso.}} Solo los tres integrantes del equipo y el
profesor titular del curso.

\textbf{\textcolor{uxnavy}{Conservación.}} El material se conserva hasta la conclusión del
curso, en agosto de 2026, y se elimina después. Las grabaciones, si las hubiera, se eliminan
una vez transcritas.

\textbf{\textcolor{uxnavy}{Participación voluntaria.}} Tu participación es voluntaria. Puedes
omitir cualquier pregunta, pedir que se detenga la sesión o retirar tu participación en
cualquier momento y sin explicación; en ese caso el material recopilado se elimina.

\textbf{\textcolor{uxnavy}{Registro del consentimiento.}} En la encuesta, el consentimiento se
registra mediante una casilla de confirmación obligatoria antes de la primera pregunta. En la
entrevista, se registra de forma verbal al inicio de la sesión, dejando constancia en las
notas.
\end{uxdestacado}

\subsection{Confirmación}

\begin{uxnota}[Declaración del participante]
He leído la información anterior, comprendo en qué consiste mi participación y acepto
participar de forma voluntaria. Entiendo que puedo retirarme en cualquier momento.

\vspace{6pt}
{\Large$\square$}\quad Acepto participar
\end{uxnota}

